%%%%%%%%%%%%%%%%%%%%%%%%%%%%%%%%%%%%%%%%%%%%%%%%%%%%%%%%%%%%%%%%%%%%%%%%%%%%%%%%%%%%%%
%%%%%%%%%%%%%%%%%%%%  INGRESAR LAS PREGUNTAS DE LA  EVALUACIÓN %%%%%%%%%%%%%%%%%%%%%%%%
%%%%%%%%%%%%%%%%%%%%%%%%%%%%%%%%%%%%%%%%%%%%%%%%%%%%%%%%%%%%%%%%%%%%%%%%%%%%%%%%%%%%%%%

\paragraph{I} Limits. Properties and Algebraically. 

\begin{preguntas}

\pregunta[6] Compute the limit of the function \(f(x) = \frac{x^2 - 4}{x - 2}\) as $x\to2$. 

\begin{center}
    \framebox(\textwidth,150){}   
\end{center}

\pregunta[7] Compute the limit of the function \(f(x) = \frac{6x^2 + 4}{5x^2 - 2}\) as $x\to\infty$.

\begin{center}
    \framebox(\textwidth,150){} 
\end{center}

\newpage

\pregunta[7] Compute the limit of the function \(f(x) = \frac{\sqrt{x+4} - 5}{x}\) as $x\to0$.

\begin{center}
    \framebox(\textwidth,150){} 
\end{center}

\end{preguntas}


\paragraph{II} Derivative using the power rule and properties of derivatives.

\begin{preguntas}

\pregunta[10]{Using the power rule, compute the derivative of the function \(f(x) = \dfrac{3}{2x}\).}

\begin{center}
    \framebox(\textwidth,150){} 
\end{center}


\pregunta[10]{Using the power rule, compute the derivative of the function \(g(x) = 4\left[\left(\dfrac{x}{2}\right)^{12} - \left(\dfrac{x}{2}\right)^{6}\right] \).}

\begin{center}
    \framebox(\textwidth,150){} 
\end{center}

\end{preguntas}


\newpage

\paragraph{III} Derivatives of trascendental functions.

\begin{preguntas}

\pregunta[6] Applying the rules of derivatives for trascendental functions, compute the derivative of the function $f(x) = \sin(x) + \cos(x)$. 

\begin{center}
    \framebox(\textwidth,150){} 
\end{center}


\pregunta[6] Applying the rules of derivatives for trascendental functions, compute the derivative of the function $f(x) = e^x + \ln\qty(x)$.

\begin{center}
    \framebox(\textwidth,150){} 
\end{center}

\pregunta[8] Applying the rules of derivatives for trascendental functions, compute the derivative of the function $f(x) = \ln\qty(5x)$.

\begin{center}
    \framebox(\textwidth,150){} 
\end{center}



\end{preguntas}

\newpage

\paragraph{IV} Derivative by definition. \(\dv{f(x)}{x} = \lim_{h\to0}\frac{f(x+h) - f(x)}{h}\).

\begin{preguntas}

\pregunta[20] Compute the derivative of the function \(f(x) = \sqrt{x}\) using the definition of derivatives. 

\begin{center}
    \framebox(\textwidth,200){} 
\end{center}

\pregunta[20] Compute the derivative of the function \( f(x) = |x| \) at $x=0$ using the definition of derivatives.

\begin{center}
    \framebox(\textwidth,200){} 
\end{center}

\end{preguntas}

